\section{Технология виртуальных полигонов}
\label{ch1_1_1}

Специфика исследовательского проектирования морских объектов состоит в неопределенности и неполноте информации об их поведении в экстремальных ситуациях, совокупно вызванных воздействием внешней среды, инициируемых им опасных ситуаций (например, интенсивным обледенением) и внутренних дестабилизирующих факторов (смещения груза, затопления отсеков и пр.). При этом развитие аварийной ситуации происходит по различным сценариям, которые не могут быть в общем случае сведены к ограниченному набору типовых случаев. Однако ретроспективный анализ гибели судов позволяет выделить условия, способствующие возникновению аварийных ситуаций. При движении судна лагом к волне выделяют такие ситуации, как воздействие на судно ветрового шквала в условиях сильной качки, потеря остойчивости по причине затопления палубного колодца, а также ударное воздействие гребня разрушающейся волны. Напротив, при попутном волнении начинают играть роль совершенно другие критические факторы, а именно: параметрический резонанс бортовых колебаний, потеря управляемости на гребне волны, захват судна волной (брочинг). В свою очередь, развитие каждой из вышеперечисленных ситуаций может усложняться за счет внутренних факторов (смещение навалочного груза, затопление отсеков, интенсивное обледенение и пр.). Неоднозначность условий, приводящих к возникновению аварийных ситуаций, в данном случае затрудняет обоснование оптимальных проектных характеристик судна формальными методами и приводит к необходимости использования методов виртуального моделирования. Иными словами, основным инструментом поддержки принятия решений в процессе исследовательского проектирования становятся численные модели динамики морских объектов \citep{nechaev_ai}, позволяющие количественно воспроизвести их поведение под воздействием реалистичных внешних возмущений.

Таким образом, единственным способом моделирования рассматриваемых ситуаций является программное численное моделирование. В связи с этим вводится понятие виртуального полигона (Virtual test bed).
Виртуальный полигон --– это программно-аппаратное средство, которое поддерживает виртуальное проектирование, анализ и прототипирование различных систем \citep{vtb_ship_ee} \citep{vtb_ILRO}. 

Виртуальные полигоны --- понятие достаточно широкое, и видение архитектуры систем разработчиками, может сильно разниться. Так, например, ключевыми особенностями системы для задач анализа корабельных электрических систем \citep{vtb_ship_ee} являются: поддержка моделей, написанных на разных языках программирования; система визуализации включающая анимацию в реальном времени и построение графиков; отчуждение пользовательских знаний.

Другая система, разрабатываемая в NASA для задач запуска суборбитальных аппаратов  \citep{vtb_ILRO},
представляет собой web окружение, основной задачей которого является объединений в общей сложности 44 категорий моделей, каждая из которых реализует свои отдельные, сравнительно простые задачи в единую распределенную систему.



%Виртуальный полигон должен удовлетворять следующим требованиями:
%\begin{itemize}
%	\item 	Реальное время выполнения (в идеале --- «сверхреальное»). «Сверхреальное» 
%			время выполнения позволит перебирать несколько вариантов параллельно, что может быть полезно для предсказания поведения реальных объектов, а также производить массовую статистическую выборку.
%	\item 	Интерактивность. В данном случае под интерактивностью подразумевается 
%			возможность изменения параметров моделируемых объектов и получения 
%			ответной реакции в реальном времени.
%	\item 	Гибкость. Виртуальный полигон должен предоставлять набор конструктивных 
%			и вспомогательных объектов для формирования задачи моделирования.  
%			При этом объекты могут быть как конструктивными, так и вспомогательными. 
%			Конструктивные отвечают за определение и формирование поведения и свойств 
%			реального объекта, а вспомогательные позволяет формировать сценарии 
%			модельных процессов и обработки результатов.
%	% \item Возможность самонастройки и самоподстройки. 
%	% Виртуальный полигон должен обеспечивать возможность гибкого 
%	\item 	Возможность расширения: возможность добавления новых явлений, свойств и объектов.
%\end{itemize}

Анализ задач моделирования динамики судна \citep{nechaev_ai}, а также виртуальных полигонов, использующихся в других научных и инженерных областях показал, что основной частью виртуального полигона является интегрирующая различные модели подсистема (интегратор). Также в состав виртуального полигона может входить некоторый 
\frqt{конструктор} для формирования объектов и средство анализа. Конструирование объектов может осуществляться как посредством ввода команд, задающих параметры новых объектов так и с использованием развитых WYSIWYG-редакторов. Средство анализа результатов моделирования может представлять как простой построитель графиков по данным, так и сложные системы 3D-визуализации. Компоненты виртуального полигона могут быть как локальными (собранными в одно приложение и развернутым на ПК оператора), так и распределенными (развернутыми на нескольких серверах, со множеством клиентов).

Обобщенная схема виртуального полигона продемонстрирована на рис.~\ref{vtb_concept}

\begin{figure}
\begin{center}
	\includegraphics[width=130mm]{vtb_concept}
	\label{vtb_concept}
	\caption{Обобщенная схема виртуального полигона}
\end{center}
\end{figure}

%\begin{itemize}
%	\item   Подсистема человеко-машинного взаимодействия, в состав которой входят:
%		\begin{itemize}
%		\item 	Подсистема визуализации и отображения данных, 
%				которая отвечает за отображение текстовой информации, 
%				построения диаграмм а также двумерной, трехмерной и 
%				стерео-визуализации моделируемого процесса.
%		\item 	Подсистема ввода, которая может включать разнообразные 
%				устройства ввода, начиная от простейшего текстового 
%				интерфейса терминального типа и заканчивая интерфейсами 
%				основанными на захвате движения.
%		\end{itemize}
%	\item	Подсистема моделирования --- является центральной 
%			частью виртуального полигона и отвечает за расчет параметров 
%			моделируемого объекта в каждый момент времени. Подсистема 
%			моделирования может быть как монолитной, так и композитной. 
%			В последнем случае, каждая компонента подсистемы моделирования 
%			выполняет свою задачу и обменивается данными с другими компонентами.
%	\item	Подсистема сценариев --- подсистема, которая отвечает за 
%			формирование последовательностей, условий и дополнительных 
%			действий, выполняемых в процессе моделирования.
%\end{itemize}
%
%Подсистемы виртуального полигона могут быть либо собраны в один програмный пакет, запускаемый на персональной ЭВМ так и сформированны в распределенный программно-аппаратный комплекс, который может состоять из:
%\begin{itemize}
%	\item 	Вычислительных серверов, кластеров, на которых 
%			могут быть развернуты компоненты подсистемы моделирования.
%	\item 	Графических станции, которые формируют изображение 
%			целиком или по фрагментами. Разные графические могут 
%			визуализировать моделируемый процесс с разных точек, 
%			в разных масштабах и режимах отображения.
%	\item 	Компонент ввода, которые могут включать в себя клавиатуры, 
%			мыши, сенсорные панели, системы захвата движения, нейро-интерфейсы и т.п.
%\end{itemize}

Как было уже отмечено выше --- центральной частью виртуального полигона является 
подсистема моделирования (локальная или распределенная). В основе подсистем моделирования 
применительно к задачам экстремальной динамики судна лежат модели внешней среды 
(нерегулярное морское волнение) и модель поведения судна. 

Выбор комплекса моделей с одной стороны должен отвечать требованиям задачи, которую необходимо решить, 
с другой стороны требованиями по вычислительным трудозатратам. В свою очередь комплекс моделей определяет
способы взаимодействия между ними, а также форматы входных и выходных данных.
